\section{Conclusions}
\label{sec:conclusions}

This chapter presented DDMORL, a decay-driven multi-objective reinforcement learning framework for single-node energy management in HPC environments. The central contribution is a formulation that lets users specify their acceptable performance degradation as a single scalar $\delta \in [0,1]$, which is automatically converted to a preference-aligned control signal via a lightweight application-specific progress-to-PCAP model $\mathcal{G}_{app}$. This eliminates the need for abstract preference weight tuning, learned preference interpolators, or separate agents per operating point, all of which are practical barriers to deploying MORL in production HPC settings.

\noindent\textbf{Summary of contributions.}
The framework trains a single Q-network offline on pre-collected execution traces from a mixed dataset of eight benchmarks, comprising five STREAM memory-bandwidth variants and three NAS Parallel Benchmark variants. At deployment, the controller receives the user-specified $\delta$, computes the aligned preference vector $\boldsymbol{\omega}$, and selects PCAP actions using a preference-aware scoring function that jointly maximises alignment and expected return. No online exploration, no system downtime for retraining, and no enumeration of static power caps are required.

\noindent\textbf{Experimental findings.}
Evaluation on nine benchmarks, comprising all five STREAM variants and four NAS Parallel Benchmarks, confirms three key properties of the framework. First, the preference knob $\delta$ reliably steers the controller across the full energy-performance trade-off surface: a single trained controller covers all 14 discrete power-cap levels in the action space $\{89, 95, \ldots, 165\}$\,W, with PCAP selection shifting monotonically from the high end (approximately $130$--$165$\,W) for $\delta \leq 0.2$ to the low end (approximately $89$--$107$\,W) for $\delta \geq 0.7$. Second, the STREAM benchmarks achieve energy savings of $31$--$39\%$ with well-separated Pareto operating points across the full $3$--$95\%$ degradation range, while the NPB benchmarks achieve $20$--$31\%$ savings, with operating-point clustering in the high-degradation regime reflecting the saturation behavior of compute-bound workloads. Third, the overlay of DDMORL operating points against static fixed-PCAP baselines in the execution-time versus energy plane shows that the controller's output is indistinguishable from that of an optimally tuned static power cap, providing empirical evidence that the policy inference step introduces no measurable runtime or energy overhead.

\noindent\textbf{Generalisation.}
Trained on a mixed dataset of five STREAM variants and three NPB variants, with NPB-FT held out as an out-of-distribution test case, the controller produces coherent and monotone Pareto fronts across all nine evaluated benchmarks. The successful transfer to NPB-FT is attributed to the PAPI-derived state features (IPC, L3 miss rate, and stall ratio), which encode workload-agnostic microarchitectural characteristics that are predictive of the application's response to PCAP changes. The result demonstrates that DDMORL learns a preference-conditioned power-management strategy rather than a workload-specific lookup table.

\noindent\textbf{Limitations and future directions.}
The current framework operates at the single-node level and assumes that progress is approximately monotonic with respect to PCAP, a property that holds for the memory-bandwidth-bound and embarrassingly parallel workloads studied here but may not hold for tightly coupled, communication-intensive applications at scale. Extending DDMORL to multi-node scenarios would require coordinating PCAP decisions across nodes while accounting for inter-node communication bottlenecks, which introduces a substantially more complex state and action space. Additionally, the $\mathcal{G}_{app}$ model is constructed from training-time data for a fixed application; adapting it online to new workloads or hardware configurations without full retraining is an open problem. Future work will explore hierarchical control architectures that combine the node-level DDMORL controller developed in this chapter with the job-scheduler-level policies presented in earlier chapters, enabling end-to-end energy optimisation from the scheduling layer down to the hardware actuator.
